(1) 选择向右为正，滑块坐标为 \((X+s\cos\theta,-s\sin\theta)\)。由此 \(L=\frac12(M+m)\dot X^2+m\dot X\dot s\cos\theta+\frac12m\dot s^2+mg s\sin\theta\)。约束是几何性质而非另外施加的虚拟力；\(X\) 不显含于拉氏量，是循环坐标。
(2) \(p_X=(M+m)\dot X+m\dot s\cos\theta=0\)，因为初态静止且不存在外界水平冲量。于是 \(\dot X=-m\dot s\cos\theta/(M+m)\)，代回得到约化质量 \(\mu=m(M+m\sin^2\theta)/(M+m)\) 和 \(\mu\ddot s=mg\sin\theta\)，即 \(\ddot s=(M+m)g\sin\theta/(M+m\sin^2\theta)\)。
(3) 将上述 \(\ddot s\) 代入动量关系，可得 \(\ddot X=-mg\sin\theta\cos\theta/(M+m\sin^2\theta)\)。取斜面向外法向 \(\mathbf n=(\sin\theta,\cos\theta)\)，牛顿第二定律在该方向给 \(N-mg\cos\theta=m\ddot X\sin\theta\)，故 \(N=mMg\cos\theta/(M+m\sin^2\theta)>0\)，在题设正质量、正角度下始终维持接触。
(4) 消去水平坐标后 \(E=\frac12\mu\dot s^2-mgs\sin\theta\) 为常量，其导数等于 \(\dot s(\mu\ddot s-mg\sin\theta)=0\)。固定楔块极限 \(M\to\infty\) 给 \(\ddot s=g\sin\theta\)、\(N=mg\cos\theta\)；水平面极限 \(\theta\to0\) 给加速度趋零、\(N\to mg\)，均与直观一致。
